%============================================================ 7. 고찰
\section{고찰}
\label{sec:discussion}

\subsection{계획 안정성이라는 설계 원리}

\S\ref{sec:res_mechanism}의 관찰 --- 커버리지는 성능을 예측하지 못하고 churn과 경로 교차가
예측한다 --- 은 본 문제에 국한되지 않는 함의를 갖는다.

주기적으로 재계획하는 계층 구조에서는 상위 계층의 결정 변경이 하위 계층에 \emph{전환 비용}을
발생시킨다. 본 문제에서 그 비용은 구체적이다. 방어정이 그물벽을 절반쯤 전개한 상태에서 담당이
바뀌면, 그때까지 소모한 이동과 그물이 회수되지 않는다. 상위 계층이 매 재계획마다 순간
최적을 추구하면 이 비용이 누적되어, 국소적으로 더 나은 배정이 전역적으로 손해가 된다.

이는 계획 주체를 언어모델로 바꾼다고 자동으로 해결되지 않는다. 오히려 언어모델은 매 호출이
독립적이므로 \emph{일관성이 구조적으로 보장되지 않는다}. 본 연구에서 직전 배정을 상태에
포함시키고 담당 무리를 방위로 식별하게 한 것은 이 때문이다. 무리 식별자는 매 결정마다 방위
정렬 순서로 다시 매겨지므로, 식별자만 보면 같은 무리인데도 담당이 사라진 것처럼 보여 불필요한
전환이 발생한다.

따라서 언어모델을 계획 계층에 도입할 때의 설계 지침은 다음과 같이 정리된다. 계획 품질을
\emph{커버리지나 국소 최적성}으로 평가하지 말고 \emph{시간에 걸친 일관성}으로 평가해야 하며,
프롬프트와 상태 표현이 그 일관성을 지원해야 한다.

\subsection{조합 최적화를 제거한 선택의 대가}

\S\ref{sec:commander}에서 배정 후처리를 제거한 것은 측정 가능성을 위한 선택이었다. 그 대가가
\S\ref{sec:res_threshold}의 음의 결과로 나타난다. 후처리가 있었다면 Qwen2.5 7B의 배정 오류는
상당 부분 코드에 의해 교정되어 $-0.192$ 라는 악화가 관측되지 않았을 것이다.

그러나 그 경우 얻어지는 것은 ``LLM 지휘관이 잘한다''가 아니라 ``후처리가 잘한다''이다.
본 연구의 예비 측정에서 후처리를 둔 구성은 배정안의 64.5\,\%를 변경하였고 14개 상황 중
57\,\%에서 언어모델 출력과 무관하게 동일한 최종 배정에 수렴하였다. 2-opt 개선이 사실상
완전탐색에 가까워 초기값이 무엇이든 같은 국소최적으로 수렴하기 때문이다.

이 점은 언어모델 기반 계획 연구 전반에 대한 방법론적 지적이기도 하다. 언어모델의 출력을
최적화나 규칙으로 보정한 뒤 종단 성능만 보고하면, 보고된 성능이 언어모델에서 온 것인지
보정기에서 온 것인지 원리적으로 구분되지 않는다. 기여를 주장하려면 보정을 제거한 조건이
함께 보고되어야 한다.

\subsection{로컬 대 클라우드 배치}

\S\ref{sec:res_latency}의 지연 측정은 배치 선택에 직접적인 함의를 갖는다. 로컬 14B는
클라우드 대비 10.6배 느리면서 성능도 낮았다(0.811 대 0.904). 해상 방어에서 오프라인 운용이
요구조건이라면, 현시점에서 온보드 언어모델 지휘관은 다음 두 조건을 동시에 만족해야 한다.

첫째, \S\ref{sec:res_threshold}의 능력 임계를 넘어야 한다. 본 실험에서 14B는 넘지 못했다.
둘째, 지연이 결정 주기 안에 \emph{여유를 남기고} 들어와야 한다. 14B의 평균 16.59초는 25초
주기 안에 들어오지만, 실측에서 호출의 4.2\,\%가 주기를 초과했고 최댓값은 296초였다
(Table~\ref{tab:latency_tail}). 동기 구조라면 그 4.2\,\%마다 제어 루프가 정지한다. 비동기
설계는 이 꼬리 위험을 제거한다는 점에서 로컬 배치의 전제 조건에 가깝다.

이는 평균 지연만 보고 배치를 판단하면 안 된다는 실무적 함의를 갖는다. 다만 이 지연은 시뮬레이터와 GPU 를 공유한 구성에서 측정한 값이므로, 절대값은 \S\ref{sec:limitations} 의 교란 요인과 함께 읽어야 한다. 평균으로는 주기의
3분의 2에 불과해 여유가 있어 보이지만, 분포의 최댓값은 주기의 12배에 달한다.

다만 이 결론은 시점 의존적이다. 동일 파라미터 규모의 모델 품질은 계속 향상되고 있으므로,
임계를 넘는 로컬 모델의 크기는 낮아질 것으로 예상된다. 본 연구가 제시하는 것은 특정 모델의
순위가 아니라 \emph{임계가 존재하며 그것을 측정하는 방법}이다.

\subsection{점수맵 표현의 이득}

\S\ref{sec:res_effectsize}에서 가장 큰 개선이 나타난 지표가 총 선회량($\dz = -3.23$)이라는
점은 점수맵 표현의 성질과 직접 연결된다. 경유점이 격자 위에 고정되므로 연속 회귀에서
발생하는 미세 진동이 구조적으로 배제되고, 그 결과 궤적이 부드러워진다. 선회 감소는 단순한
심미적 개선이 아니라 실제 선박에서 연료·시간·기동 여유로 환산되는 양이다.

포획당 그물 소모의 감소($-45\,\%$) 역시 표현과 연관된다. 유효 마스크가 배정 회랑과 사거리
제약을 행동 분포에 직접 부과하므로, 애초에 무의미한 위치에 그물벽을 놓는 행동이 표현되지
않는다. 손실 항으로 제약을 유도하는 방식에서는 이런 행동이 낮은 확률로나마 남는다.
